% =============================================================================
%  Tiểu kết chương 2
% =============================================================================

\section*{Tiểu kết}
\addcontentsline{toc}{section}{Tiểu kết}
Chương này đã trình bày các cơ sở lý thuyết và các hướng nghiên cứu liên quan đến đề tài, bao gồm các phương pháp biểu diễn tri thức bằng Ontology và đồ thị tri thức, các hệ thống hỏi đáp dựa trên đồ thị tri thức, cũng như kiến trúc Retrieval-Augmented Generation và các mở rộng như GraphRAG. Các phương pháp này cho thấy tiềm năng lớn trong việc kết hợp tri thức có cấu trúc với các mô hình sinh ngôn ngữ nhằm nâng cao chất lượng truy vấn và trả lời câu hỏi.

Tuy nhiên, phần lớn các nghiên cứu hiện nay chủ yếu tập trung vào các nguồn dữ liệu quy mô lớn hoặc các miền tri thức phổ quát, trong khi việc xây dựng và khai thác đồ thị tri thức cho các miền tri thức văn hóa chuyên biệt vẫn còn hạn chế. Từ những phân tích trên, khóa luận xây dựng đồ thị tri thức cho miền tri thức nghệ thuật Chèo, đồng thời thiết kế hệ thống GraphRAG nhằm khai thác tri thức này trong các tác vụ truy vấn và hỏi đáp. Các phương pháp và kiến trúc của hệ thống sẽ được trình bày chi tiết trong chương tiếp theo.